\section{Experimental Setup}
\label{sec:setup}

\subsection{Datasets and Flow Partitions}
Experiments are conducted across two prominent network intrusion detection benchmarks:
\begin{enumerate}
    \item \textbf{Source Domain: CIC-IDS-2017~\cite{sharafaldin2018cicids}:} Generated by the Canadian Institute for Cybersecurity, this dataset captures benign activity and twelve distinct attack profiles (including Brute Force, DoS, Heartbleed, Web Attacks, and DDoS) recorded over five consecutive days. The raw corpus contains $2,830,743$ flows. We extract a stratified, balanced working dataset of $358,018$ flows, partitioned into an $80\%$ training split ($N_{\text{train}} = 286,414$) and a $20\%$ in-domain holdout evaluation split ($N_{\text{test}} = 71,604$). The test split contains $56,000$ benign flows and $15,604$ attack flows ($21.8\%$ positive prevalence).
    \item \textbf{Target Domain: UNSW-NB15~\cite{moustafa2015unsw}:} Generated by the Australian Centre for Cyber Security using the IXIA PerfectStorm tool, this target corpus contains realistic modern background traffic and nine synthetic attack families. We evaluate the entire standardized testing partition comprising $N_{\text{target}} = 257,673$ flows, consisting of $93,000$ benign flows and $164,673$ malicious flows ($63.9\%$ positive prevalence). Table~\ref{tab:attack_distribution} summarizes the distribution across the nine attack categories.
\end{enumerate}

\begin{table}[h]
\centering
\caption{UNSW-NB15 Attack Category Breakdown ($N_{\text{target}} = 257,673$)}
\label{tab:attack_distribution}
\renewcommand{\arraystretch}{1.1}
\begin{tabular}{lrr}
\hline
\textbf{Category} & \textbf{Flow Count ($N$)} & \textbf{Prevalence (\%)} \\
\hline
Normal (Benign) & 93,000 & 36.09\% \\
Generic & 58,871 & 22.85\% \\
Exploits & 44,525 & 17.28\% \\
Fuzzers & 24,246 & 9.41\% \\
DoS & 16,353 & 6.35\% \\
Reconnaissance & 13,987 & 5.43\% \\
Analysis & 2,677 & 1.04\% \\
Backdoor & 2,329 & 0.90\% \\
Shellcode & 1,511 & 0.59\% \\
Worms & 174 & 0.07\% \\
\hline
\textbf{Total} & \textbf{257,673} & \textbf{100.00\%} \\
\hline
\end{tabular}
\end{table}

\subsection{Hyperparameters and Training Configuration}
The 3-layer MLP is trained using pure NumPy mini-batch Adam optimization with learning rate $\alpha = 0.001$, momentum parameters $\beta_1 = 0.9$, $\beta_2 = 0.999$, and numerical stabilizer $\epsilon = 10^{-8}$. Mini-batch size is set to $256$ flows across $20$ epochs with early stopping triggered if holdout validation loss fails to improve over three consecutive epochs. Model weights are initialized using He (Kaiming) normal initialization. To verify that observed effects are not artifacts of lucky initializations, five independent training runs are executed across random seeds $\mathcal{S} = \{42, 101, 202, 303, 404\}$.

\subsection{Evaluation Metrics and Statistical Tests}
To capture performance under extreme class imbalance, we report:
\begin{itemize}
    \item \textbf{Accuracy (Acc):} $(TP + TN) / N$.
    \item \textbf{Balanced Accuracy (BAcc):} $\frac{1}{2}\left(\frac{TP}{TP + FN} + \frac{TN}{TN + FP}\right)$, weighting benign and malicious recall equally.
    \item \textbf{Precision, Recall (TPR), F1-Score, and FPR:} Standard diagnostic measures with decision threshold $\tau = 0.5$.
    \item \textbf{McNemar's Test for Paired Classifications~\cite{mcnemar1947note}:} Given paired model predictions $A$ (FP32) and $B$ (INT8) on identical flows, the contingency table defines discordant pairs: $b = \sum \mathbb{I}(A \text{ correct } \land B \text{ incorrect})$ and $c = \sum \mathbb{I}(A \text{ incorrect } \land B \text{ correct})$. Under the null hypothesis $H_0$ that quantization does not alter classification accuracy, the Edwards continuity-corrected statistic:
    \begin{equation}
    \chi^2 = \frac{(|b - c| - 1)^2}{b + c} \sim \chi^2_1.
    \end{equation}
    \item \textbf{Bootstrap Confidence Intervals~\cite{efron1994introduction}:} We compute empirical $95\%$ confidence intervals on the quantization mitigation delta $\Delta_{\text{BAcc}} = \text{BAcc}_{\text{INT8}} - \text{BAcc}_{\text{FP32}}$ across $B = 1,000$ bootstrap resamples of the target dataset.
    \item \textbf{Wilcoxon Signed-Rank Test~\cite{wilcoxon1945individual}:} Evaluates whether continuous predicted probability distributions exhibit non-zero median shifts between FP32 and INT8 representations.
\end{itemize}

\subsection{Benchmarking Hardware Environment}
Host CPU execution timings are captured on an AMD64 Windows 11 platform equipped with a 12-core/16-thread Intel processor, running 64-bit Python 3.11.9 and NumPy 2.4.6. To eliminate OS scheduler jitter, execution benchmarks execute $200$ discarded warm-up loops, followed by $3,000$ individually timed inferences across $5$ repeats with garbage collection disabled (\texttt{gc.disable()}).
